% BHU PYQ -- Professional Exam Template
\documentclass[11pt,a4paper]{article}

\usepackage[a4paper,top=2.5cm,bottom=2.5cm,left=2.8cm,right=2.8cm,headheight=14pt]{geometry}
\usepackage{amsmath,amssymb}
\usepackage[T1]{fontenc}
\usepackage[utf8]{inputenc}
\usepackage{lmodern}
\usepackage{microtype}
\usepackage{fancyhdr}
\usepackage{enumitem}
\usepackage{listings}
\usepackage{hyperref}
\usepackage{lastpage}
\usepackage{parskip}

\hypersetup{colorlinks=true,urlcolor=black,linkcolor=black,
  pdftitle={BVCA-304: Java Applet Programming -- BHU PYQ}}

\pagestyle{fancy}
\fancyhf{}
\renewcommand{\headrulewidth}{0.4pt}
\renewcommand{\footrulewidth}{0.4pt}
\fancyhead[L]{\small\textit{Banaras Hindu University}}
\fancyhead[R]{\small\textit{BVCA-304: Java Applet Programming}}
\fancyfoot[C]{\small Page \thepage\ of \pageref{LastPage}}

\lstset{
  basicstyle=\small\ttfamily,
  frame=single,
  breaklines=true,
  columns=flexible,
  keepspaces=true
}

\newlist{parts}{enumerate}{2}
\setlist[parts,1]{label=(\alph*),leftmargin=*,itemsep=4pt,topsep=3pt}
\setlist[parts,2]{label=(\roman*),leftmargin=*,itemsep=2pt,topsep=2pt}
\newcommand{\pts}[1]{\hfill\textit{[#1]}}

\begin{document}

\begin{center}
  {\large\bfseries BANARAS HINDU UNIVERSITY}\\[2pt]
  {\normalsize B.Voc. Semester III Examination, 2023-24}\\[6pt]
  \rule{0.8\linewidth}{0.5pt}\\[6pt]
  {\large\bfseries Paper: BVCA-304 --- Java Applet Programming}\\[4pt]
  {\normalsize Time: Three Hours \hfill Maximum Marks: 70}
\end{center}

\rule{\linewidth}{0.4pt}

\smallskip
\noindent\textit{Instructions: Attempt all questions. Symbols have their usual meanings.}

\bigskip

BVCA-304 : Java Applet Programming
Answer All Sections.
Candidates are required to write their answers in their own words as far as practicable.

\bigskip\noindent{\large\bfseries SECTION --- A}
\rule{\linewidth}{0.4pt}\smallskip

Long Answer type questions to be answered in about 500 words each: 12x2=24 Marks
Attempt any two:

\medskip\noindent\textbf{Question 1.} What do you understand by the concept of Object Oriented Language? Explain how Java is
specified as an Object Oriented Language.

\medskip\noindent\textbf{Question 2.} Explain about for, while and do-while loop structures in Java with example. Why they are
\_ categorized as entry-controlled and exit controlled loops?

\medskip\noindent\textbf{Question 3.} What are operators ? Explain all the Operators supported by Java programming language.
4, Discuss Switch Statement. Write a program to calculate percentage of total marks,
obtained in 5 subjects by a student and display the equivalent grade as S, A,B,C, D, or F.
[91-100\% as S grade; 81-90\% as A grade; 71-80\% as B grade and so on. F grade represents
| Fail]
, Section - B
5, Attempt any six (Answer in about 200 words each): 6 x 6 = 36 Marks
\begin{parts}
  \item What do you understand by classes and objects? How objects are created from a class?
  \item What are arrays? How will you implement one dimensional and two dimensional arrays?
  \item Define JDK and JVM. Why Java is considered as platform independent language?
  \item ' Differentiate between static, instance and local variables. :
  \item ‘Discuss break and continue statement in java. .
  \item Write a program Count the digits of a number given input by a user.
  \item ) Write a java program find the area and circumference of a circle using the concept of
  objects and classes.
  \item ' Discuss any four inbuilt methods for String manipulation given in java library.
\end{parts}

\bigskip\noindent{\large\bfseries SECTION --- C}
\rule{\linewidth}{0.4pt}\smallskip


\medskip\noindent\textbf{Question 6.} Attempt all the questions: 10 x 1=10 Marks
\begin{parts}
  \item What does an Eclipse IDE allow?
  g a) Editing the program
  \item Compiling the program .
  \item Both Editing and Compiling the program
  \item Editing, Building, Debugging and Compiling the program.
\end{parts}

\medskip\noindent\textbf{Question 9.} 7-0.
Si abla
s
  \begin{parts}
    \item Which keyword is used for accessing the features of a package?
  \end{parts}
\begin{parts}
  \item package
  \item import
  \item extends
  \item export
  \begin{parts}
    \item | Which of these values is returned by read() method is end of file (EOF) is encountered? ,
  \end{parts}
  \item 0
  \item 1
  \item -1
  \item Null
  \begin{parts}
    \item Which one of the following is not a Java feature?
  \end{parts}
  \item Object-oriented
  \item Use of pointers
  F c) Portable
  \item Dynamic and Extensible :
  \begin{parts}
    \item JREstandsfor\_\_.
  \end{parts}
  \item Java run ecosystem
  \item JDK runtime Environment
  \item Java Runtime Environment
  \item None of these
  \begin{parts}
    \item Which class in Java is used to take input from the user?
  \end{parts}
  \item Scanner
  \item Input
  \item Applier
  \item None of these
  \begin{parts}
    \item "Automatic type conversion is possible in which of the possible cases?
  \end{parts}
  \item Byte to int
  \item Int to long
  \item Long to int
  \item Short to int
  \begin{parts}
    \item Identify the output of the following program.
    String str = “abcde”;
    System.out. printin(str.substring(1, 3));
  \end{parts}
  \item abc
  Coundi-   G)
  @
  \item bc
  \item bed
  \item cd
  \begin{parts}
    \item Find the output of the following code.
    int ++a = 100;
  \end{parts}
\begin{lstlisting}[language=C]
System.out.printin(++a);
\end{lstlisting}
  \item 101
  \item Compile error as ++a is not valid identifier
  \item 100
  \item None
  i (X) \_ Find the output of the following code.
  Public class Solution{
\begin{lstlisting}[language=C]
Public static void main(String args[]){ .
Int i;
for(i = 1; i < 6; i++){
| if(i > 3) continue;
| }
| System.out.print!n(i);
| }
\end{lstlisting}
  \item 3
  \item 4
  5
  \item 6
  oS
\end{parts}

\vfill
\rule{\linewidth}{0.4pt}
\begin{center}\small
\href{https://bhu-exam.vercel.app/}{Click here to visit our website}\\[4pt]\href{https://me-aryan.vercel.app/}{Click here to visit our contributor}\\[4pt]\href{https://quashan-me.vercel.app/}{Click here to visit our contributor}
\end{center}

\end{document}
